\section{The question and its context}\label{sec:introduction}

At aspect ratio $\lambda=3$, a numerical search finds a stationary radius near $0.823938$ in reference coordinates. Rotating that point gives a circle. What excludes a second radius, or a stationary point outside the equatorial plane? An exploratory three-dimensional cover of the normalized radial derivative $H$ left 21,254 of 31,340 terminal boxes without a separated sign. Those boxes count neither stationary points nor failures of a classification theorem. They tell us exactly which question the first experiment could not answer.

The second experiment used the vertical derivative away from the equator, then counted scalar crossings on the equator. A third experiment, reported here, asks whether the stronger equatorial monotonicity suggested by point values survives a whole-domain enclosure. The first experiment covers a solid meridian domain and a shape interval; the third covers the equatorial radius and shape only. We use each experiment on its own domain. Together they motivate the question for the full family: how do the number, dimension and type of stationary components change with shape?

An extremizing centre is one part of this question. Saddles also matter, because they record changes in the stationary structure which an extremum selector would discard. Spheroids provide a particularly accessible test: the sphere has one stationary point, but rotational symmetry permits both a pair of isolated axial critical points and a circular critical manifold away from the centre. The purpose of the present classification is to determine when these sets occur and to exclude additional components over a noncompact shape interval.

\input{sections/related_work}

The proof has three levels. First, analytic formulas identify the scalar derivatives whose signs exclude off-axis or off-equatorial stationary points. Second, one-variable arguments determine the number and type of the remaining zeros. Third, finite interval covers and explicit estimates at the two degenerating ends provide the signs used in these arguments. The label \emph{basepoint geometry} refers here to the first formulation, and \emph{GeoEnclosure} to this application of established enclosure methods to the resulting expressions. Neither label is a claim that interval arithmetic, cone volume measures, or critical manifolds originate in this work.

\section{Energy, coordinates and boundary convention}\label{sec:definition}

\subsection{A moving-basepoint energy}
Let $K\subset\R^3$ be a smooth strictly convex body, let $p$ be an interior point, and let $n(x)$ denote the outward unit normal at $x\in\partial K$. Define
\begin{align}
 u_p(x)&=\frac{x-p}{|x-p|},&
 \alpha_p(x)&=\arccos\bigl(u_p(x)\cdot n(x)\bigr),\label{eq:angle}\\
 d\mu_{K,p}(x)&=\frac{(x-p)\cdot n(x)}{3\vol(K)}\,dA(x),&
 E_K(p)&=\int_{\partial K}\alpha_p(x)^2\,d\mu_{K,p}(x).\label{eq:energy}
\end{align}
Angles are measured in radians. Strict interiority implies $(x-p)\cdot n(x)>0$, and the divergence theorem gives
\[
 \int_{\partial K}(x-p)\cdot n(x)\,dA(x)=3\vol(K).
\]
Thus $\mu_{K,p}$ is a probability measure. Both the angle and the weight depend on $p$ and must be differentiated. Pushing this boundary measure forward by the normal map gives the cone volume measure of $K-p$, normalized by volume \citep{Boroczky2013,Huang2016}.

The energy is invariant under translations, orthogonal changes of coordinates and positive similarities. For a scale factor $a$, the numerator, surface element and volume in \eqref{eq:energy} scale by $a,a^2,a^3$, respectively. An arbitrary affine map changes angles and is not an invariance of this functional.

There is a geometric interpretation of the integrand. Write the boundary as $x=p+r_p(u)u$, $u\in S^2$. Its outward normal is parallel to $u-\nabla_{S^2}\log r_p$. Radial projection, or taking the scalar product of the oriented area element with $u$, gives $(x-p)\cdot n\,dA=r_p^3\,d\omega$. Hence
\[
 \tan\alpha_p=|\nabla_{S^2}\log r_p|,\qquad
 E_K(p)=\frac{1}{3\vol K}\int_{S^2}r_p^3
 \arctan^2\!\bigl(|\nabla_{S^2}\log r_p|\bigr)\,d\omega .
\]
The angle measures the directional variation of log radial distance. The energy vanishes exactly when the radial distance is constant, that is, when $K$ is a ball centred at $p$.

\subsection{Volume preserving spheroids}
Write
\begin{equation}\label{eq:family}
 K_\lambda=M_\lambda\B,\qquad
 M_\lambda=\diag(\lambda^{-1/3},\lambda^{-1/3},\lambda^{2/3}),\qquad\lambda>0.
\end{equation}
Here $\B$ is the closed unit ball, and its interior is understood whenever a basepoint is specified. The determinant is one and $\vol(K_\lambda)=4\pi/3$. The shapes are oblate for $\lambda<1$, spherical for $\lambda=1$, and prolate for $\lambda>1$. Neither $\lambda=0$ nor $\lambda=\infty$ is included as a member of the family.

Set $p=M_\lambda q$, $|q|<1$, and $\E_\lambda(q)=E_{K_\lambda}(M_\lambda q)$. Similarity invariance allows us to use $L_\lambda=\diag(1,1,\lambda)$. For $\xi\in\Sph$, put
\begin{align}
 k&=1-q\cdot\xi,& D&=|L_\lambda(\xi-q)|,& S&=|L_\lambda^{-\mathsf T}\xi|,\label{eq:reference_variables}\\
 \gamma&=\frac{k}{DS},& f(c)&=\arccos^2c,&
 \av{h}&=\frac1{4\pi}\int_{\Sph}h(\xi)\,d\Omega.\nonumber
\end{align}
Then
\begin{equation}\label{eq:reference_energy}
 \E_\lambda(q)=\av{kf(\gamma)}.
\end{equation}
Indeed, for $x=M\xi$, the surface Jacobian is $\det M\,|M^{-\mathsf T}\xi|$ and the normal is $M^{-\mathsf T}\xi/|M^{-\mathsf T}\xi|$. Multiplying the Jacobian by $(x-p)\cdot n=k/|M^{-\mathsf T}\xi|$ and using $3\vol(M\B)=4\pi\det M$ proves \eqref{eq:reference_energy}.

For fixed $\lambda$ the chain rule gives
\begin{equation}\label{eq:coordinate_hessian}
 \nabla_q\E_\lambda=M_\lambda^{\mathsf T}\nabla_pE_{K_\lambda},\qquad
 \nabla_q^2\E_\lambda=M_\lambda^{\mathsf T}(\nabla_p^2E_{K_\lambda})M_\lambda.
\end{equation}
Zeros and Hessian inertia are preserved. Numerical derivative bounds throughout this manuscript use the reference variable $q$; their magnitudes must not be copied unchanged to the physical coordinates.

Rotations about the third axis and reflection across the equator reduce the energy to $\mathcal E_\lambda(\rho,z)=\E_\lambda(\rho,0,z)$, with $\rho\ge0$. It is even in $z$. The centre is stationary for every $\lambda>0$. Interior compact subsets have $D$ bounded away from zero. The apparent singularity of derivatives of $f$ at $c=1$ is removable, so the interior differentiations used below are legitimate.

\subsection{A bounded first-derivative kernel}
Define the continuously extended functions
\begin{align}
 P(c)&=-f'(c)=\frac{2\arccos c}{\sqrt{1-c^2}},& Q(c)&=2P(c)+cP'(c),\label{eq:PQ}\\
 J(c)&=cP(c)-f(c),& K(c)&=c^2P(c).\nonumber
\end{align}
For $0\le c\le1$,
\begin{equation}\label{eq:PQbounds}
 P(c)=2\int_0^\infty\frac{dt}{\cosh t+c},\quad
 2\le P\le\pi,\quad -2\le P'\le0,\quad Q\ge\frac{10}{3}.
\end{equation}
The first identity follows by $s=\tanh(t/2)$. Differentiation under its integral sign yields the bounds on $P'$ and
\[
 Q(c)=2\int_0^\infty\frac{2\cosh t+c}{(\cosh t+c)^2}\,dt.
\]
The integrand decreases with $c$, while $P(1)=2$, $P'(1)=-2/3$ and $Q(1)=10/3$.

With $u=L_\lambda(\xi-q)/D$, differentiation of \eqref{eq:reference_energy}, using $k_i=-\xi_i$ and $D_i=-(L_\lambda)_{ii}u_i$, gives
\begin{equation}\label{eq:gradient}
 \partial_{q_i}\E_\lambda
 =\av{\xi_iJ(\gamma)-(L_\lambda)_{ii}S u_iK(\gamma)}.
\end{equation}
For each fixed finite $\lambda>0$, the kernel is uniformly bounded for interior $q$. At a boundary limit $q\to q_0\in\Sph$, it converges pointwise except at the single contact point $\xi=q_0$, a set of sphere measure zero. Dominated convergence gives a unique continuous trace $g_\partial(q_0;\lambda)$ of the full gradient. The domination is also uniform when $\lambda$ ranges over a compact subset of $(0,\infty)$.

\begin{definition}[Boundary zero convention]\label{def:boundary}
The boundary zero set is
\begin{equation}\label{eq:boundary}
 Z_\partial(\lambda)=\{q\in\Sph:g_\partial(q;\lambda)=0\}.
\end{equation}
It is the zero set of the full gradient trace, not the constrained stationarity condition $(I-qq^{\mathsf T})g_\partial=0$ on the sphere.
\end{definition}

Only first-derivative traces are used. In particular, no continuous extension of the boundary Hessian is assumed. A signed axial component at a general boundary point establishes a nonzero gradient component; it does not by itself determine the sign of the outward normal derivative.

\section{Classification and its explicit inputs}\label{sec:classification}

Let $h_c(\lambda)=\partial_{q_1q_1}\E_\lambda(0)$ and $a_c(\lambda)=\partial_{q_3q_3}\E_\lambda(0)$. The four critical values are denoted by
\[
 0<\cO<\bO<1<\bP<\cP<\infty.
\]
They are specified by the zeros and rational intervals in Table~\ref{tab:critical}. The table uses endpoint signs for each defining function. The earlier local enclosure for the prolate centre is compared below. Decimal approximations in the final column indicate location.

\begin{table}[htbp]\centering\small
\caption{Critical parameters. $(a,b)/2^m$ means $(a/2^m,b/2^m)$. Decimal values indicate location only.}\label{tab:critical}
\begin{tabular}{@{}llll@{}}\toprule
Value & Defining zero & Recorded rational interval & Reference decimal\\\midrule
$\cO$ & $a_c$ & $(407958860300,407958860301)/10^{12}$ & $0.40795886030$\\
$\bO$ & polar axial trace & $(674802,674811)/2^{20}$ & $0.64354577$\\
$\bP$ & equatorial radial trace & $(34651360,34651369)/2^{24}$ & $2.06538229$\\
$\cP$ & $h_c$ & $(79261996,79262005)/2^{24}$ & $4.72438340$\\\bottomrule
\end{tabular}
\end{table}

Write $O=\{0\}$ and, when they exist, let
\begin{align}
 A_\lambda&=\{(0,0,z_*(\lambda)),(0,0,-z_*(\lambda))\},\quad 0<z_*<1,\label{eq:axial_set}\\
 C_\lambda&=\{(\rho_*(\lambda)\cos\phi,\rho_*(\lambda)\sin\phi,0):0\le\phi<2\pi\},\quad 0<\rho_*<1.\label{eq:circle_set}
\end{align}
Let $P_\pm=\{(0,0,1),(0,0,-1)\}$ and let $C_\partial$ be the unit equatorial circle. Physical sets are obtained by applying $M_\lambda$.

\begin{theorem}[Classification from Assertions C1--C6]\label{thm:classification}
Assume the explicitly stated computational Assertions C1--C6 in Section~\ref{sec:computational}. For the energy \eqref{eq:reference_energy}, the interior stationary set is $O$ together with the set in the third column of Table~\ref{tab:classification}, and the full-gradient boundary zero set is the fourth column. There are no additional stationary components.
\end{theorem}

\begin{table}[htbp]\centering\small
\caption{Classification in reference coordinates. Index counts negative Hessian eigenvalues; for a circle it is the normal index.}\label{tab:classification}
\begin{tabularx}{\linewidth}{@{}lXXX@{}}\toprule
Shape range & Centre & Other interior set & Boundary zeros\\\midrule
$(0,\cO)$ & nondegenerate saddle, \mbox{index 1} & none & none\\
$\cO$ & degenerate saddle, nullity 1 & none & none\\
$(\cO,\bO)$ & nondegenerate minimum & $A_\lambda$: two index 1 saddles & none\\
$\bO$ & nondegenerate minimum & none & $P_\pm$\\
$(\bO,\bP)$ & nondegenerate minimum & none & none\\
$\bP$ & nondegenerate minimum & none & $C_\partial$\\
$(\bP,\cP)$ & nondegenerate minimum & $C_\lambda$: Morse--Bott saddle circle & none\\
$\cP$ & degenerate saddle, nullity 2 & none & none\\
$(\cP,\infty)$ & nondegenerate saddle, \mbox{index 2} & none & none\\\bottomrule
\end{tabularx}
\end{table}

The interval $(\bO,\bP)$ contains the sphere. A circle is one connected critical manifold, not a finite collection of isolated roots. Its boundary arrival at $\bP$, and the axial pair's arrival at $\bO$, are excluded from the interior count. This classification does not select all constrained global minimizers on the closed body: a constrained boundary minimum can have nonzero normal derivative.

The theorem is stated with a named, finite list of inputs so that its current proof status can be assessed. In a certified submission version, C1--C6 must be supplied as established computational propositions with an accessible, fixed reproduction package. Removing the assumption from the wording without establishing those propositions would not strengthen the result.

\noindent\textbf{Scope of the separate oblate verification.}
The author's 16 November 2026 record
(\href{https://github.com/cotaxxxx/bg-spheroid-family/blob/127be52046e9498d1306fe3d35c3e02e91c636de/runs/2026-11-16/D_OB_GLOBAL_CLOSED_V24/JUDGE_RULING_D_OB_GLOBAL_CLOSED_20261116.md}{fixed record}; see Section~\ref{sec:availability})
covers only $\lambda\in[2/5,33/50]$.
A separate project AI audit checked the assembly and source citations,
not the computations anew. Neither record establishes C1--C6 or
covers all oblate shapes.

\section{Two sign principles and the sphere}\label{sec:principles}

\subsection{Symmetry factors and crossings}
Use the continuous extensions
\begin{equation}\label{eq:HV}
 H(\rho,z;\lambda)=\frac{\partial_\rho\mathcal E_\lambda}{\rho},\qquad
 V(\rho,z;\lambda)=\frac{\partial_z\mathcal E_\lambda}{z}.
\end{equation}
For example,
\[
 V(\rho,z;\lambda)=\int_0^1\partial_{zz}\mathcal E_\lambda(\rho,tz)\,dt.
\]
A nonzero $H$ excludes a stationary point only when $\rho>0$; a nonzero $V$ does so only when $z\ne0$.

\begin{lemma}[One-way crossings]\label{lem:crossing}
Let $h$ be $C^1$ on an interval. If $h'(x)<0$ at every zero, then $h$ has at most one zero.
\end{lemma}
\begin{proof}
Each zero is isolated. In a compact interval between two hypothetical zeros, an infinite zero sequence would accumulate at a zero with derivative zero; hence there are finitely many. Two consecutive zeros are impossible: immediately after the first the sign is negative, and immediately before the second it is positive, forcing another zero in between.
\end{proof}

The sign rule has a degree-theoretic counterpart. Let $U\subset\mathbb R^n$ be bounded and open, $F$ continuous on $\overline U$ and $C^1$ on $U$, with $0\notin F(\partial U)$ and $0$ a regular value. The preimages are isolated and form a compact set, hence are finite. Their signed count is $\deg(F,U,0)=\sum_{F(x)=0}\operatorname{sgn}\det DF(x)$. When the signs agree, the unsigned count is the absolute degree. This observation does not determine the degree without boundary information; our scalar proof uses the two endpoint signs directly.

For a box $B$, target domain $\mathcal D$ and scalar quantity $F$, an enclosure means
\begin{equation}\label{eq:enclosure}
 \{F(q,\lambda):(q,\lambda)\in B\cap\mathcal D\}\subseteq I_F(B).
\end{equation}
Exclusion of zero proves a sign. Inclusion of zero does not prove the existence of a zero. Existence and uniqueness below use continuity, endpoint signs and Lemma~\ref{lem:crossing}; separate components containing zero need not vanish at a common point.

\subsection{The spherical case, proved in the text}\label{sec:sphere}
For $\lambda=1$, rotational invariance gives $\E_1(q)=e(|q|)$. Put $q=re_3$ and $\xi\cdot e_3=\cos\psi$. With
\[
 \theta(r,\psi)=\arctan\frac{r\sin\psi}{1-r\cos\psi},\qquad
 \Delta=1-2r\cos\psi+r^2,
\]
we have
\begin{equation}\label{eq:sphere_energy}
 e(r)=\frac12\int_0^\pi(1-r\cos\psi)\theta(r,\psi)^2\sin\psi\,d\psi.
\end{equation}

\begin{proposition}\label{prop:sphere}
The sphere has exactly one interior stationary point, its centre. It is the unique global minimum and its Hessian is $(4/3)I$. The boundary zero set in Definition~\ref{def:boundary} is empty.
\end{proposition}
\begin{proof}
Using $\theta_r=\sin\psi/\Delta$ and $\theta_\psi=r(\cos\psi-r)/\Delta$, differentiation followed by integration by parts yields
\begin{equation}\label{eq:sphere_derivative}
 e'(r)=\frac12\int_0^\pi
 \frac{\theta(r,\psi)\sin^2\psi}{\Delta}
 (2-r^2-r\cos\psi)\,d\psi>0\qquad(0<r<1).
\end{equation}
The last factor is at least $(1-r)(r+2)>0$, and the other factors are positive in the interior of the integration interval. Since $\nabla\E_1(q)=e'(|q|)q/|q|$, only the centre is stationary. There $e(0)=0$, whereas the integrand in \eqref{eq:sphere_energy} is positive on a set of positive measure for every nonzero basepoint.

For $|r|<1$, the angle also satisfies
$\theta(-r,\psi)=-\theta(r,\pi-\psi)$.
Substituting $\psi\mapsto\pi-\psi$ in \eqref{eq:sphere_energy}
therefore yields $e(-r)=e(r)$.
Hence every odd Taylor coefficient of $e$ vanishes.
The absolutely convergent expansion
$\theta=\sum_{n\ge1}r^n\sin(n\psi)/n$,
uniformly convergent on compact intervals of $|r|<1$,
permits termwise integration and gives
\begin{align}
 e(r)&=\frac23r^2-\frac29r^4-\frac{22}{1575}r^6+O(r^8),\label{eq:sphere_series}\\
 \frac{e'(r)}r&=\frac43-\frac89r^2-\frac{44}{525}r^4+O(r^6).\nonumber
\end{align}
This proves the centre Hessian. To justify taking the boundary limit
in \eqref{eq:sphere_derivative}, observe that for $1/2\le r<1$,
\[
 \Delta=(1-r)^2+2r(1-\cos\psi)\ge2r(1-\cos\psi),
 \qquad
 \frac{\sin^2\psi}{\Delta}
 \le\frac{1+\cos\psi}{2r}\le\frac1r.
\]
Also $0\le\theta(r,\psi)\le\pi/2$ and
$|2-r^2-r\cos\psi|\le3$ in this range.
Thus the integrand in \eqref{eq:sphere_derivative} is bounded by a
constant independent of $r$ and $\psi$; dominated convergence applies.
For $0<\psi<\pi$,
$\theta(1,\psi)=(\pi-\psi)/2$ and
$\Delta(1,\psi)=2(1-\cos\psi)$. Consequently,
\[
 e'(1^-)=\frac14\int_0^\pi\frac{\pi-\psi}{2}\sin^2\psi\,d\psi
 =\frac{\pi^2}{32}>0.
\]
By rotational symmetry the resulting full-gradient trace is
$g_\partial(q_0;1)=e'(1^-)q_0$ for every $q_0\in\Sph$.
This is exactly the trace in Definition~\ref{def:boundary},
not a constrained tangential gradient on the boundary.
It is nonzero everywhere on $\Sph$.
\end{proof}

The two terms $4/3-8r^2/9$ are a Taylor approximation, not an exact identity or a global lower bound. None of the finite-shape lower bounds used later is derived from that truncation. Proposition~\ref{prop:sphere} has no dependency on an unpublished sphere manuscript.

\section{Analytic closure of the flat end}\label{sec:oblate_analytic}

This section proves the two estimates used for $0<\lambda\le1/32$. They control the full ball, including approaches in which the basepoint and shape degenerate at unrelated rates. We use $\xi=(x,y,\mu)$, $\beta=1-\lambda^2$, $s^2=1-\mu^2$ and $m=|\mu|$.

\begin{lemma}[Equatorial reflection inequality]\label{lem:equatorial}
For $0<\lambda<1$ and $0<\rho\le1$,
\begin{equation}\label{eq:equatorial_bound}
 \partial_\rho\mathcal E_\lambda(\rho,0)\ge\frac{\lambda\rho}{24}>0.
\end{equation}
At $\rho=1$ the derivative means the interior trace.
\end{lemma}
\begin{proof}
Start with $0<\rho<1$ and put $\delta=1-\rho^2$. An integration by parts in azimuth in \eqref{eq:gradient} gives
\begin{equation}\label{eq:equator_identity}
 H(\rho,0;\lambda)=\av{\frac{P(\gamma)}{SD^3}
 \left[\delta y^2+\mu^2\{1+\beta s^2-(\rho+\beta/\rho)x\}\right]}.
\end{equation}
This identity may also be obtained by differentiating the azimuthal periodic quantity $y f(\gamma)$ and integrating its derivative to zero; the boundary terms vanish by periodicity.

Suppose first that $0<\rho\le\beta$. Set $a=\rho/\beta$ and $T=\beta/\rho-\rho\ge0$. The braces in \eqref{eq:equator_identity} equal
\[
 \beta((x-a)^2+y^2)-T(x-a).
\]
Project both hemispheres onto the unit disk $\mathbb D$. Writing $w=\sqrt{\lambda^2+\beta m^2}=\lambda S$, the signed contribution is
\begin{equation}\label{eq:reflection_I}
 I=\av{\frac{P(\gamma)\mu^2(x-a)}{SD^3}}
 =\frac{\lambda}{2\pi}\int_{\mathbb D}\frac{x-a}{D^3}\frac m w P(\gamma)\,dx\,dy.
\end{equation}
Pair a point $(a+t,y)$, $t>0$, with $(a-t,y)$. If the right point belongs to the disk, so does the left. Put $r^2=t^2+y^2$, $A=1-\rho^2/\beta>0$ and $C=A+\beta r^2$. Direct expansion gives
\begin{align}
 D_+^2=D_-^2&=\lambda^2A+\beta r^2,&
 w_\pm^2&=A-\beta r^2\mp2\rho t,\label{eq:mirror_equalities}\\
 \gamma_\pm&=\frac{\lambda}{2D}\left(w_\pm+\frac C{w_\pm}\right).\nonumber
\end{align}
We have $w_-\ge w_+$ and $w_-w_+\le C$. The difference of $w+C/w$ between the two values has the sign of $(w_--w_+)(1-C/(w_-w_+))$, so $\gamma_-\le\gamma_+$. The function $P$ decreases, whereas $m/w$ increases with $m$. The left weight $(m/w)P(\gamma)$ is consequently at least the right weight. Since $D_+=D_-$, each paired contribution to $I$ is nonpositive. Unpaired points lie to the left and also contribute nonpositively. Hence $I\le0$.

The term $-TI$ can now be discarded in a lower bound. Using $P\ge2$, $S\le1/\lambda$, $D\le2$, and
\[
 \av{y^2}=\av{\mu^2}=\frac13,\qquad
 \av{\mu^2(x^2+y^2)}=\frac2{15},\qquad \av{\mu^2x}=0,
\]
we obtain
\begin{equation}\label{eq:equator_inner}
 H\ge\frac\lambda4\left[\frac13+\frac{2\beta}{15}
       +\frac{\rho^2}{3}(\beta^{-1}-1)\right]\ge\frac\lambda{12}.
\end{equation}
For $\beta\le\rho<1$, the braces in \eqref{eq:equator_identity} are at least $(1-\rho s)(1-\beta s/\rho)\ge0$. On the hemisphere $x\le0$ they are at least one. As $\av{\mu^2\one_{x\le0}}=1/6$, this gives $H\ge\lambda/24$. Taking the first-derivative trace as $\rho\uparrow1$ proves \eqref{eq:equatorial_bound}.
\end{proof}

\begin{lemma}[Uniform axial concavity of thin oblate spheroids]\label{lem:axial_concavity}
For $0<\lambda\le1/32$ and $\rho^2+z^2<1$,
\begin{equation}\label{eq:thin_concavity}
 \partial_{zz}\mathcal E_\lambda(\rho,z)\le-\frac6{55}\lambda.
\end{equation}
\end{lemma}
\begin{proof}
Write $u_h=(x,y)$, $p_h=(\rho,0)$, $R=|u_h-p_h|$, $v=\mu-z$ and $h_1=\lambda^2kv/D^2$. Since $D^2=R^2+\lambda^2v^2$, differentiating \eqref{eq:gradient} gives
\begin{align}
 \partial_{zz}\mathcal E&=P_+-N_z,\label{eq:positive_negative}\\
 P_+&=\av{\frac{P\lambda^2k^2R^2}{SD^5}},&
 N_z&=\av{\frac Q{SD}(\mu-h_1)^2}.\nonumber
\end{align}
For example, $\gamma_z=-(\mu-h_1)/(SD)$, $(v/D)_z=-R^2/D^3$, $J'=Q$ and $K'=\gamma Q$ verify the decomposition directly.

First, the spherical average satisfies
\begin{equation}\label{eq:R_inverse}
 \av{R^{-1}}=\frac\pi2\qquad(0\le\rho<1).
\end{equation}
To see this, project onto the horizontal disk and take polar coordinates about $p_h$. On a ray of angle $\vartheta$, put $b_\vartheta=\sqrt{1-\rho^2\sin^2\vartheta}$. Then $m^2=b_\vartheta^2-(R+\rho\cos\vartheta)^2$. The radial integral $\int dR/m$ is $\pi/2-\arcsin(\rho\cos\vartheta/b_\vartheta)$. The second term cancels against the opposite ray. Both hemispheres and the normalization $1/(2\pi)$ give \eqref{eq:R_inverse}.

The identity $k=u_h\cdot(u_h-p_h)+\mu v$ yields $k^2\le2R^2+2m^2v^2$. The first term contributes at most $2\pi\lambda^2\av{R^{-1}}=\pi^2\lambda^2$ to $P_+$, by $P\le\pi$, $S\ge1$ and $D\ge R$. For the second term, use $S\ge m/\lambda$ and project the hemispheres. Its contribution is at most
\[
 \frac{\lambda^3}{2}\sum_{\sigma=\pm1}\int_{\mathbb D}
 \frac{v^2R^2}{(R^2+\lambda^2v^2)^{5/2}}\,dx\,dy,
 \qquad v=\sigma m-z.
\]
For $R\le4\lambda$, set $t=\lambda|v|/R$ and use $t^2/(1+t^2)^{5/2}\le1/4$ to bound the density by $1/(4\lambda^2R)$. For $R>4\lambda$, $|v|\le2$ bounds it by $4/R^3$. Extending each upper bound to the plane gives
\[
 2\pi\left[\int_0^{4\lambda}\frac{dR}{4\lambda^2}
       +\int_{4\lambda}^\infty\frac4{R^2}\,dR\right]=\frac{4\pi}{\lambda}.
\]
It follows that
\begin{equation}\label{eq:positive_bound}
 P_+\le(\pi^2+4\pi)\lambda^2<24\lambda^2.
\end{equation}

For the negative contribution retain only the projected domain
\[
 \Omega=\{|u_h|^2\le15/16,\ R\ge1/4\}.
\]
On it, $m\ge1/4$, $|v|\le2$ and $|k|\le R+m|v|$. Thus
\[
 \frac{|h_1|}{m}\le\lambda^2\left(\frac2{mR}+\frac4{R^2}\right)
 \le\frac{96}{1024}<\frac1{10},\qquad \frac m w\ge\frac{32}{33}.
\]
The last bound follows from $w^2/m^2\le1+1/64<(33/32)^2$. Hence $|\mu-h_1|\ge9m/10$ on both hemispheres.

Removing a disk of radius $1/4$ from a disk of squared radius $15/16$ leaves area at least $7\pi/8$. Integrating $R^2$ over the larger centred disk, using $\rho^2\le1$, yields
\[
 \int_\Omega(R^2+4\lambda^2)\,dx\,dy
 \le\pi\frac{15}{16}\left(1+\frac{15}{32}+\frac1{256}\right)<\frac{3\pi}{2}.
\]
The mean is less than $12/7<16/9$, so Jensen's inequality for $u^{-1/2}$ gives
\[
 \int_\Omega\frac{dx\,dy}{\sqrt{R^2+4\lambda^2}}\ge\frac34|\Omega|.
\]
Combining this with $Q\ge10/3$, $D\le\sqrt{R^2+4\lambda^2}$ and the projection Jacobian gives
\begin{equation}\label{eq:negative_bound}
 \frac{N_z}{\lambda}\ge\frac{10}{3}\left(\frac9{10}\right)^2
 \frac{32}{33}\frac34\frac7{16}=\frac{189}{220}.
\end{equation}
The factor $7/16$ is the lower bound for $|\Omega|/(2\pi)$. Equations \eqref{eq:positive_bound} and \eqref{eq:negative_bound}, with $\lambda\le1/32$, prove
\[
 \partial_{zz}\mathcal E\le\left(\frac{24}{32}-\frac{189}{220}\right)\lambda=-\frac6{55}\lambda.
\]
\end{proof}

By equatorial reflection, $\partial_z\mathcal E(\rho,0)=0$. Integrating \eqref{eq:thin_concavity} along a vertical segment shows $V\le-6\lambda/55$ when $z\ne0$. Thus all off-equatorial points are excluded. Lemma~\ref{lem:equatorial} excludes the remaining noncentral equatorial points. At the centre the two horizontal curvatures are positive and the axial curvature is negative, so it is a nondegenerate index 1 saddle. The same quantitative first-derivative inequalities pass to the boundary trace: $z\ne0$ is excluded by the axial component, and the boundary equator by the radial component. No comparison between the rates of $\lambda\downarrow0$, $\rho\uparrow1$ and $z\to0$ is needed.

\section{Finite shapes: proof of the classification from C1--C4}\label{sec:finite_assembly}

\subsection{Oblate reduction to the axis}
Assertion C1 excludes every off-axis point for $1/32\le\lambda<1$ by $H>0$ or by a direct positive radial derivative. Put $g(z,\lambda)=\partial_z\mathcal E_\lambda(0,z)$. It is odd in $z$. The outer-axis signs in C1 and the intermediate census in C2 give
\begin{equation}\label{eq:axis_count}
 \#\{z\in(0,1):g(z,\lambda)=0\}
 =\begin{cases}
 0,&1/32\le\lambda\le\cO,\\
 1,&\cO<\lambda<\bO,\\
 0,&\bO\le\lambda<1.
 \end{cases}
\end{equation}
The existing positive root is simple with $g_z<0$, while the transverse curvatures are positive. Reflection produces two index 1 saddles. The zero at $z=0$ is never counted in \eqref{eq:axis_count}, and a root at $z=1$ is a boundary event. An intermediate cut at $z=1/2$ used in assembling the covers is a bookkeeping section, not an additional bifurcation.

For the polar trace $B_o(\lambda)=g(1,\lambda)$, the recorded endpoint bounds are
\begin{align}\label{eq:oblate_endpoint}
 B_o(674802/2^{20})&\le-175746/2^{36}<0,\\
 B_o(674811/2^{20})&\ge153930/2^{36}>0.\nonumber
\end{align}
Their signs prove existence. Assertion C2 supplies uniqueness in the remaining polar window, and C1 excludes further polar zeros outside it. At $\lambda=\bO$, only the two poles belong to the boundary zero set. Off-axis boundary exclusion uses the directly enclosed trace, or the quantitative bounds specified in C1; mere strict positivity without a uniform margin would not justify passage to a boundary limit.

The endpoint behaviour shown in Figure~\ref{fig:branches} also follows
conditionally from these same gates.  As $\lambda\downarrow\cO$ along the
unique positive-root branch, any accumulation point of $z_*(\lambda)$
in $(0,1]$ would be a positive zero of $g(\cdot,\cO)$, including its
polar trace at $1$; C2 and the negative polar sign exclude this.
Thus $z_*(\cO+)=0$.  As $\lambda\uparrow\bO$, C2 excludes every
positive interior zero at $\bO$.  Moreover
$g_z(0,\bO)=a_c(\bO)>0$, so by continuity $g(z,\lambda)/z>0$
uniformly for sufficiently small $z>0$ and $\lambda$ close to $\bO$.
The branch cannot return to the centre, and its only possible
accumulation point is $1$. Hence $z_*(\bO-)=1$.
This argument uses the C2 root exclusions and trace continuity; it
is not a new numerical certification.

Near the oblate centre, write $g=a_c(\lambda)z+c_3(\lambda)z^3+O(z^5)$. C2 gives $a_c(\cO)=0$, $a_c'(\cO)>0$ and $c_3(\cO)<0$. At $\cO$ the axial energy restriction has negative leading term $c_3(\cO)z^4/4$, while horizontal restrictions have positive quadratic terms. This proves that the degenerate centre is a saddle of nullity one.

\subsection{Near-sphere prolate band}
For $1\le\lambda\le2$, Assertion C3 provides
\begin{equation}\label{eq:near_V}
 V(\rho,z;\lambda)\ge16189/2^{36}>0.
\end{equation}
The bound is obtained from a direct $E_z$ enclosure away from the equator and from an $E_{zz}$ enclosure integrated from the equator. It includes $\lambda=1$ directly and does not extrapolate a truncated spherical expansion.

On $z=0$, the additional equatorial input in C3 gives $H_\rho=\rho E_{yyyy}/3<0$ for $0<\rho<1$, by the rotation identity derived below. Its positive boundary trace therefore implies $H(\rho,0;\lambda)>0$ for every interior radius, and also at the centre by continuity. Together with the positive axial curvature, this makes the centre a nondegenerate minimum.

\subsection{The prolate window and its circle}
On $2\le\lambda\le5$, C4 gives
\begin{equation}\label{eq:window_V}
 V(\rho,z;\lambda)\ge3341881/2^{36}>0.
\end{equation}
Again only the equatorial plane remains. Let $h_c(\lambda)=H(0,0;\lambda)$ and $h_b(\lambda)=H(1^-,0;\lambda)$. The unique zeros of these functions are $\cP$ and $\bP$, respectively.

At $(q_1,q_2,q_3)=(\rho,0,0)$, $\rho>0$, rotational symmetry gives
\begin{equation}\label{eq:fourth_derivative}
 \partial_{q_1q_1}\E=H+\rho H_\rho,\qquad
 \partial_{q_2q_2q_2q_2}\E=\frac{3H_\rho}{\rho}.
\end{equation}
For the second identity, expand $\widehat E_\lambda(\sqrt{\rho^2+y^2},0)$ at $y=0$. The whole-equator inequality in C4 now gives $H_\rho=\rho E_{yyyy}/3<0$ throughout $0<\rho<1$. At fixed shape, continuity and strict decrease give exactly one interior radius when $h_c>0>h_b$. The endpoint signs therefore select precisely $\bP<\lambda<\cP$. If either endpoint vanishes, strict decrease excludes another zero in the interior. The two integral representations used to establish this inequality appear in Section~\ref{sec:radial-experiment}.

The additional sign cover in C4 gives $H_\lambda<0$ on the zero set.
This parameter-direction sign is used to establish the decreasing
\emph{radius} of the existing saddle circle as $\lambda$ varies; the
number, existence range and normal type of circles instead follow from
vertical exclusion, $H_\rho<0$, the two endpoint signs and the normal
Hessian inequalities. The implicit function theorem then yields
\begin{equation}\label{eq:branch_slope}
 \rho_*'(\lambda)=-H_\lambda/H_\rho<0.
\end{equation}
Compactness, root uniqueness and the endpoint exclusions imply $\rho_*(\bP+)=1$ and $\rho_*(\cP-)=0$. At a circle point the normal Hessian is diagonal by reflection, with
\[
 \mathcal E_{\rho\rho}=\rho_*H_\rho<0,\qquad \mathcal E_{zz}>0,
 \qquad \mathcal E_{\rho z}=0.
\]
The tangent direction has zero curvature. Its tangent line is exactly the Hessian kernel, so the circle is Morse--Bott with normal index one \citep{Bott1954}.

Write $H(\rho,0;\lambda)=h_c(\lambda)+d_2(\lambda)\rho^2+O(\rho^4)$ near the prolate centre. C4 supplies $h_c(\cP)=0$, $h_c'(\cP)<0$ and $d_2(\cP)<0$. At $\cP$ the equatorial restriction has negative leading energy term $d_2(\cP)\rho^4/4$, whereas the axial quadratic term is positive. Hence the centre is a degenerate saddle of nullity two, not a nondegenerate index 2 saddle at that parameter.

\input{sections/radial_experiment}

\subsection{Boundary and coverage connections}
For $1\le\lambda\le5$, the uniform positive bounds in \eqref{eq:near_V} and \eqref{eq:window_V} exclude boundary points with $z\ne0$. On the boundary equator only the radial component remains; its coefficient is $h_b$, which vanishes only at $\bP$. The large-prolate argument in the next section excludes every boundary point for $\lambda\ge5$. The thin-oblate argument shares $\lambda=1/32$ with C1--C2, and the three prolate intervals share the endpoints $2$ and $5$. Together with Proposition~\ref{prop:sphere}, these overlapping ranges establish all rows of Table~\ref{tab:classification} from C1--C6.
